\section{A typical width for a linear cut budget}
\label{sec:typical_tail_scales}

The polynomial error asserted after the typical-set construction in
\cite[Section~7]{OpenAI2026AreaLaw} follows from the exponential tail estimate
of Lemma~3.1 once the cut budget is at most linear in the scale.

\begin{theorem}[Uniform polynomial decay of the tail expression]
\label{thm:typical_tail_linear_budget}
\lean{Entropy.exists_typical_tail_bound_of_linear_budget}
\leanok
Fix $\theta\ge0$ and $C>0$. There is an integer $N\ge2$, depending only on
these two constants, such that for every integer $n\ge N$ and every
$0<\mathcal B\le Cn$,
\[
2e^{e/2}\exp\!\left(-\frac{n^{3/5}}
 {32\sqrt{(1+\theta)\mathcal B}}\right)\le n^{-100}.
\]
\end{theorem}
\begin{proof}
\leanok
The denominator is at most
$32\sqrt{(1+\theta)C}\,n^{1/2}$. Consequently the left-hand side is bounded
by $2e^{e/2}\exp(-cn^{1/10})$, where
$c=(32\sqrt{(1+\theta)C})^{-1}>0$.
Exponential decay dominates every fixed real power; in particular this
expression is at most $n^{-100}$ for all sufficiently large $n$.
The threshold does not depend on the choice of $\mathcal B$.
\end{proof}

This is a numerical implication. Its physical application still requires
both the geometric bound $\mathcal B\le Cn$ and the marginal tail estimate
for the actual ground state. Given those estimates, the normalized-state
results above give positive selected mass at least $1-n^{-100}$ at width
$n^{3/5}$.
